\section{Two questions to ask the curve}\label{cq:section}

The density has already given us a way to recognize individual primes. It is
natural to ask whether familiar conjectures can be read from the probability
law without extracting its Taylor coefficients one at a time. Two observations
do this. One changes the way we sample the law and watches its variance. The
other compares two moments and watches the signs of a matrix's eigenvalues.

These are criteria: the relevant sign inequality and the unbounded matrix
index remain to be proved. Their interest is that the observables are ordinary
probabilistic and linear-algebraic objects belonging to the same law.

\subsection{The Riemann hypothesis as a variance comparison}

Write \(Y=\log W\). For \(t>-1\), change its distribution by giving an outcome
of size \(Y=y\) a weight proportional to \(e^{ty}\). More formally,
\[
 \E_t h(Y)=\frac{\E[W^t h(\log W)]}{\E W^t}.
\]
This is exponential tilting. Increasing \(t\) favors the larger values of
\(W\). If \(\kappa(t)=\log \E W^t=-\log K(t)\), differentiation gives
\begin{equation}\label{cq:variance}
 \Var_t(Y)=\kappa''(t)=\sum_p\frac{1}{(p-1+t)^2}.
\end{equation}
The variance is a softened count of primes. The prime number theorem suggests
replacing the prime atoms by the continuous density \(du/\log u\), giving
\[
 V_0(t)=\int_2^\infty\frac{du}{(u-1+t)^2\log u}.
\]
Does the actual variance eventually lie above or below this smooth comparison?

\begin{theorem}[An eventual variance deficit]\label{cq:rhvariance}
The Riemann hypothesis holds if and only if
\[
             \Var_t(\log W)<V_0(t)
             \qquad\text{for all sufficiently large }t.
\]
\end{theorem}

The word \emph{eventually} matters. This is an inequality along the positive
real axis, with no assumed simplicity or spacing of zeta zeros. Eventual-sign
criteria for RH have a long history; Robin's work \cite{Robin1984} is an
important example. The statistic here is the tilted variance of the law
already constructed.

\paragraph{What would an experimenter see?}
The parameter $t$ changes which outcomes matter most: increasing it
favors larger values of $W$. We then measure the spread of their
logarithms. The comparison curve $V_0(t)$ is obtained by replacing the
discrete primes with their smooth predicted density. The criterion
says that RH is exactly an eventual ordering of these two spreads.
No individual zero of zeta appears in the statement.

The proof exposes why a one-sided comparison is possible. Prime squares
give a negative correction to the smooth approximation. On RH, that
correction is larger than all the possible oscillations from the
nontrivial zeros combined. If RH fails, oscillations force crossings
in both directions arbitrarily far out. The complete argument and its
quantifiers are in Appendix~\ref{app:rh-variance}.

This is an exact reformulation attached to a probability law we already
met. Its usefulness as a route to RH would require a new way to compare
those variances. The theorem supplies the comparison to aim for.

\subsection{Twin primes hidden in matrices with positive entries}

For this observation we separate the prime \(2\), writing
\[
 W=2U_2V,\qquad
 K_o(s)=\prod_{p>2}(1-p^{-1})^s\left(1+\frac{s}{p-1}\right).
\]
Then \(\E V^s=1/K_o(s)\) for \(\Re s>-2\). Compare moments three units apart:
\[
 F(t)=\frac{\E V^t}{\E V^{t+3}}=\frac{K_o(t+3)}{K_o(t)},\qquad t>0.
\]
The number three has a purpose. Differences between odd primes are even;
a shift of three crosses precisely the gap of two before it reaches the next
possible gap.

For a symmetric matrix, its \emph{negative index} is the number of negative
eigenvalues, counted with multiplicity. Equivalently, it is the greatest
dimension of a subspace on which its quadratic form is strictly negative.

\begin{theorem}[A matrix count of twin pairs]\label{cq:twins}
Fix \(a>0\), and form
\[
 \mathsf H_N(a)=
 \left[\frac{(-1)^{i+j}F^{(i+j+2)}(a)}{(i+j+2)!}\right]_{0\leq i,j<N}.
\]
Every entry is strictly positive. Nevertheless,
\[
 \sup_{N\geq1}n_-(\mathsf H_N(a))
   =\#\{p>2:p\text{ and }p-2\text{ are prime}\}.
\]
Both sides are allowed to be infinite. The negative index is nondecreasing
in \(N\).
\end{theorem}

The surprise is the exact arithmetic count. A quotient of moments
of the probability law gives a sequence of positive-entry matrices whose
negative directions count twin pairs.

The general relation between signed moments and negative squares belongs to
indefinite moment theory \cite{BergChristensenMaserick1988}. We prove the
compact signed-measure step directly in the appendix, which also shows
exactly where the primes enter.

\paragraph{Why a negative direction means a twin pair.}
The moment quotient has one pole associated with each odd prime $p$.
Its residue is negative precisely when $p-2$ is also prime. This sign
test is elementary: the shift by three crosses the preceding odd prime
only when the two primes are twins.

The derivatives assemble those signed residues into a moment matrix.
Polynomials can separate any finite collection of its negative atoms,
turning each chosen twin pair into a negative direction in a sufficiently
large matrix. Conversely, the matrix cannot have more negative directions
than there are negative atoms. Appendix~\ref{app:twin-index} makes this
argument precise and controls the infinite tail.

Thus the matrices provide a counter whose reading can rise as their size
increases. The twin-prime conjecture says that this counter never reaches
a final value. All its entries are positive throughout; the signs being
counted belong to quadratic forms, which combine entries using both
positive and negative weights.

For comparison, the positive-entry matrix
$\left[\begin{smallmatrix}1&2\\2&1\end{smallmatrix}\right]$
has quadratic form $-2$ at the vector $(1,-1)$, and therefore one
negative direction. Our matrices turn this familiar linear-algebraic
possibility into an exact count of prime pairs.

\subsection{What a finite calculation can certify}

The theorem turns any rigorously negative \(k\)-dimensional subspace into a
certificate of \(k\) twin pairs. One can build the matrices from ordinary
reciprocal-prime sums. At \(a=1\), write
\[
 \frac{F(1-z)}{F(1)}=\exp\left(\sum_{k\geq1}\frac{d_kz^k}{k}\right)
                    =\sum_{n\geq0}c_nz^n,\qquad
 d_k=\sum_{p>2}\{p^{-k}-(p+3)^{-k}\}.
\]
Then
\[
 c_0=1,\qquad c_n=\frac1n\sum_{k=1}^n d_kc_{n-k},\qquad
 \mathsf H_N(1)=F(1)[c_{i+j+2}]_{0\leq i,j<N}.
\]
For an integer cutoff \(Y\),
\[
       0\leq d_k-d_k^{(Y)}\leq3Y^{-k}.
\]
These are explicit bounds for every uncomputed prime, not a decision to
ignore the tail.

An exact rational certificate supplied with the manuscript verifies that
\([c_{i+j+2}]_{0\leq i,j<15}\) has at least four negative directions using
only primes through \(100\) and the displayed tail bounds. We give its
specification in Appendix~\ref{cq:certificate}. It recovers already known
pairs; it demonstrates a verification mechanism rather than a new
twin-prime lower bound.

There is a reason to resist turning this calculation into a promise of an
easy proof of infinitude. The tail bounds allow all omitted \(d_k\)-increments
to be zero. Any conclusion valid uniformly for every such tail must also
hold for the finite prime product, whose negative index counts only its
existing twins. A proof forcing further negative directions needs additional
arithmetic information. The matrix formulation tells us precisely what that
information would have to accomplish.

The next question is different. Instead of asking a probability law to
answer a conjecture, we can ask whether its observable relationships reveal
the hidden integers that produced them.
